\section{Methodology}\label{methodology}

The UCI Diabetes 130-US hospitals (1999–2008) dataset is a publicly available dataset. It contains features that are numerical variables as well as categorical variables. Additionally, 
there are many columns with missing data as well as duplicate rows. Logistic regression expects numerical values as input features. As a result it 
is important to clean the data before it can be used.

\subsection{Data Cleaning}\label{data_cleaning}

Firstly, duplicate rows and empty columns were handled. The dataset was ordered in ascending order using the 'encounter\_id' column. Thereafter, the smallest 'encounter\_id' for a specific 'patient\_nbr'
was removed. This was also done with the intention to use the 'encounter\_id' column to perform time-aware splitting during the experiment as mentioned in Section \ref{splitting}. Thereafter, missing values were handled. 
It was observed that the 'weight' column was largely unpopulated. As a result this column was dropped. The next column with the biggest amount of missing values was the
'payer\_code' column. The 'payer\_code' column was dropped as logically this column has no statistical importance on determining whether a patient readmits or not.
Thereafter, the last three columns with missing values were dealt with: 'race', 'medical\_speciality', 'diag\_1', 'diag\_2' and 'diag\_3'. These three column required some form of
imputation as all three columns were statistically important. The imputation for each of these columns were built into a custom DataCleaner class that is called in the pipeline setup
so that when the imputation is performed, only data within a specific train, validation or test set was used. This was done to prevent data leakage. Since the 'race' column only contained
a small portion of missing values, these values were imputed using the highest probability of the 'race' column within that split. Next, missing values in the 'medical speciality' column
were classified as 'missing'. This column had over 70 categories, therefore, only the top 10 categories were selected, with the rest of data termed as 'other'. Lastly, the 'diag\_1', 'diag\_2' 
and 'diag\_3' column's missing values were imputed using the most common category in each column. In addition, there were over 700 categories in each column which was then converted using frequency
encoding.

After dealing with duplicate rows and missing values, categorical variables that did not need one-hot encoding were handled. Firstly, the 'age' column was converted to numerical values 
by taking the mean of the age range. Next, columns with single-valued categories were dropped. Thereafter, binary columns were converted to 0/1. Lastly, the target column 'readmitted' 
was converted to a binary column where 1 represented readmission within 30 days and 0 represented no readmission within 30 days.


\subsection{Model Creation}

After the data was cleaned, the model was created using a pipeline. In the pipeline, firstly, the columns in need of imputation were handled by the DataCleaner class as mentioned
in the Section \ref{data_cleaning}. Thereafter, a column transformer was used to normalize all numerical columns and one-hot encode all categorical columns. This baseline
model was used for the various initial tests as well as for hyperparameter tuning. 

